% Generado: no editar cifras manualmente.
\begin{table}[!htbp]
\centering
\caption{Descomposición de las ausencias y retornos posteriores}\label{tab:cp1-ausencias}
\small\setlength{\tabcolsep}{4pt}
\renewcommand{\arraystretch}{1.15}
\begin{tabular}{lrrrr}
\toprule
\textbf{Transición} & \shortstack[c]{\textbf{Con}\\\textbf{fallecimiento}} & \shortstack[c]{\textbf{Sin}\\\textbf{fallecimiento}} & \shortstack[c]{\textbf{Retorno}\\\textbf{posterior}} & \shortstack[c]{\textbf{Sin retorno}\\\textbf{registrado}} \\
\midrule
2001--2003 & 0 & 1\,213 & 402 & 811 \\
2003--2012 & 546 & 1\,508 & 603 & 1\,451 \\
2012--2015 & 2\,741 & 1\,010 & 393 & 3\,358 \\
2015--2018 & 1\,209 & 1\,736 & 1\,066 & 1\,879 \\
2018--2021 & 1\,135 & 999 & 0 & 2\,134 \\
\bottomrule
\end{tabular}
\par\medskip
\parbox{\linewidth}{\footnotesize Fuente: elaboración propia con agregados de CP1, ENASEMsimple v2 (2001--2021). Frecuencias sin ponderación. Las dos primeras columnas de conteos dividen las ausencias según tipent=5 en la ronda inmediatamente anterior. Las dos últimas dividen esas mismas ausencias según retorno posterior observado; no sumar las cuatro. Sin retorno registrado no significa pérdida definitiva; 2021 es el cierre del archivo.}
\end{table}
